% !TEX root = ../main.tex
\chapter{ケーススタディ：料金計算の仕様証明}
\label{chap:ch26_price_spec}

\begin{chapterabstract}
本章では、料金計算を題材に、仕様を Lean 4 で検査可能な形へ落とし込む方法を扱います。
税、割引、手数料、丸めは、いずれも日常的な業務ロジックです。
しかし、それらをどの順序で適用するかによって、結果は簡単に変わります。
本章では、金額を自然数で表す小さなモデルを作り、順序を入れ替えてよい処理と、入れ替えてはいけない処理を区別します。
さらに、等式で証明する性質、不等式で証明する性質、Lean だけでなく仕様書側で明示すべき性質を整理します。
\end{chapterabstract}

\begin{learninggoals}
\item 料金計算仕様を、金額、操作、順序、丸め、証明対象に分解できます。
\item 税、割引、手数料のうち、順序を入れ替えてよい処理と入れ替えてはいけない処理を区別できます。
\item Lean 4 で、価格計算の等価性、反例、不等式による安全性を表せます。
\item リファクタリング時に「証明できた範囲」と「仕様として残す範囲」を説明できます。
\end{learninggoals}

\section{この章を読む理由}

料金計算は、見た目より壊れやすい処理です。
税込金額を計算します。
割引を適用します。
送料や決済手数料を足します。
キャンペーン上限を適用します。
最後に丸めます。
これらは、単体では難しい処理に見えません。
しかし、処理を組み合わせると、順序の違いがそのまま金額の違いになります。

たとえば、次の二つは同じでしょうか。

\begin{itemize}
\item 商品価格から割引し、その後に手数料を足します。
\item 商品価格に手数料を足し、その後に割引します。
\end{itemize}

多くの場合、この二つは同じではありません。
手数料を割引対象に含めるのか、含めないのかという業務仕様の違いになるためです。
一方、すでに割引後の金額に「固定額の税」と「固定額の手数料」を足すだけなら、足す順番を変えても合計は変わりません。

図\ref{fig:ch26-price-pipeline}では、割引、手数料、税、丸めを別々の処理として並べ、順序と丸め位置を仕様の一部として表します。

\FigurePlaceholder{fig-ch26-price-pipeline.png}{0.90}{料金計算パイプライン}{fig:ch26-price-pipeline}

本章の狙いは、実際の課金システムを完全に検証することではありません。
税法、通貨単位、決済代行会社の仕様、会計上の丸め規則、キャンペーン適用条件などは、現実には大きな問題です。
ここでは、設計レビューに持ち込める小さな核を作ります。
つまり、「このリファクタリングでは計算結果が変わらない」「この順序変更は一般には安全でない」「丸めの位置は仕様として固定する必要がある」といった主張を、Lean で読める形にします。

\section{ソフトウェア上の問題}

EC、SaaS、サブスクリプション、予約システム、請求書発行などでは、料金計算の実装が少しずつ成長します。
最初は単純な合計金額だけだったとしても、後から次のような仕様が追加されます。

\begin{itemize}
\item クーポン割引を適用します。
\item 決済手数料や送料を加えます。
\item 消費税やサービス税を計算します。
\item 端数を切り捨てるか、切り上げるか、四捨五入します。
\item 割引上限、最低請求額、無料条件を入れます。
\end{itemize}

このような処理は、コード上では関数の合成として現れます。
たとえば、\texttt{discount}、\texttt{addFee}、\texttt{addTax}、\texttt{round} を順に適用するパイプラインです。

この章では、まず金額を「最小通貨単位の自然数」として表します。
日本円なら円、米ドルならセントのように、整数で扱える単位まで落とすイメージです。
小数や実数ではなく自然数から始めるのは、Lean で最小モデルを作りやすくするためです。

\begin{warningbox}
本章のモデルは、実務の金額型そのものではありません。
自然数の引き算は 0 で止まるため、過大な割引を適用しても負の金額にはなりません。
この振る舞いを採用するか、エラーにするか、符号付き整数を使うかは、実システムの仕様として別に決める必要があります。
\end{warningbox}

\section{料金仕様を証明しやすい形に分解する}

自然言語の仕様をそのまま Lean に入れることはできません。
まず、仕様を小さな部品に分けます。

\begin{definitionbox}
本章では、料金計算仕様を次の五つに分けます。
\begin{enumerate}
\item \textbf{金額の表現}：金額をどの型で表すかを決めます。
\item \textbf{操作}：割引、手数料、税、丸めをどの関数として表すかを決めます。
\item \textbf{順序}：操作をどの順に合成するかを決めます。
\item \textbf{証明対象}：等式、不等式、反例のどれとして主張するかを決めます。
\item \textbf{証明しない範囲}：外部仕様、法的ルール、小数丸めなどをどこまでモデル化しないかを決めます。
\end{enumerate}
\end{definitionbox}

この分解は、仕様レビューでも有効です。
次の表は、本章で使う対応関係です。

\begin{table}[htbp]
\centering
\caption{料金仕様の対応}
\label{tab:ch26-spec-decomposition}
\begin{tabular}{lll}
\toprule
自然言語の仕様 & Lean での表現 & 証明したい性質 \\
\midrule
金額 & \texttt{Money := Nat} & 金額を自然数として扱う \\
固定手数料 & \texttt{addFee} & 金額を減らさない \\
固定額の税 & \texttt{addFixedTax} & 手数料と順序交換できる \\
割引 & \texttt{applyDiscount} & 金額を増やさない \\
率による税 & \texttt{taxOnly} & 丸め位置に注意する \\
計算パイプライン & \texttt{totalV1}, \texttt{totalV2} & リファクタリング前後の一致 \\
\bottomrule
\end{tabular}
\end{table}

ここで大切なのは、すべてを一度に証明しようとしないことです。
まず、固定額の税や手数料のように、足し算だけで説明できる部分を扱います。
次に、割引のように順序の影響が出る部分を見ます。
最後に、率と丸めのように、仕様の境界を明示しなければならない部分を扱います。

\section{Lean 4 で最小モデルを作る}

まず、金額と基本操作を定義します。
ここでは \texttt{Money} を \texttt{Nat} の別名にします。
\texttt{Nat} は自然数なので、0 以上の整数だけを表します。

\begin{listing}[htbp]
\begin{minted}{lean4}
import Std

namespace Chapter26

abbrev Money := Nat

def applyDiscount (discount amount : Money) : Money :=
  amount - discount

def addFee (fee amount : Money) : Money :=
  amount + fee

def addFixedTax (tax amount : Money) : Money :=
  amount + tax

def taxOnly (rateNum rateDen amount : Money) : Money :=
  (amount * rateNum) / rateDen

def addRateTax (rateNum rateDen amount : Money) : Money :=
  amount + taxOnly rateNum rateDen amount
\end{minted}
\caption{\texttt{Money} の基本操作}
\label{lst:ch26-money-model}
\end{listing}

\texttt{applyDiscount} は、金額から割引額を引きます。
Lean の \texttt{Nat} では、\texttt{500 - 1000} は負の数ではなく \texttt{0} になります。
これは「割引後の金額は 0 未満にならない」という仕様に近い振る舞いですが、現実のシステムでは、過大割引をエラーにする設計もあります。

\texttt{taxOnly} は、\texttt{amount * rateNum / rateDen} によって税額を計算します。
たとえば \texttt{rateNum = 1}、\texttt{rateDen = 10} なら、10\% 相当の税額です。
ただし、\texttt{/} は自然数の整数除算なので、端数は切り捨てられます。
つまり、この関数にはすでに丸めが含まれています。

また、型は \texttt{rateDen = 0} を排除していません。

本章の反例は分母 \texttt{10} だけを使いますが、一般の税率 API にするなら、分母が正である証拠を引数または型に含める必要があります。

\section{丸め、税、割引、手数料をどう並べるか}

料金計算で最初に決めるべきことは、どの金額を対象に、どの操作を適用するかです。
たとえば、手数料に割引を適用するのでしょうか。
税込後に割引するのでしょうか、それとも割引後に税を計算するのでしょうか。
明細ごとに丸めるのでしょうか、それとも合計後に丸めるのでしょうか。

この章では、まず次の標準順序を使います。

\begin{enumerate}
\item 商品基準額 \texttt{base} から割引 \texttt{discount} を引きます。
\item 割引後の金額に固定手数料 \texttt{fee} を足します。
\item 最後に固定額の税 \texttt{tax} を足します。
\end{enumerate}

ここで「固定額の税」という単純化をしている点に注意してください。
現実の税は多くの場合、率で計算され、端数処理があります。
その部分は、後の節で明示的に扱います。

\section{順序を入れ替えてよい処理}

まず、入れ替えてよい例から見ます。
固定税と固定手数料は、どちらも単に金額を足す処理です。
そのため、同じ金額に対して両方を足すなら、順序を変えても結果は同じになります。

Lean では、この性質を次のように書けます。

\begin{listing}[htbp]
\begin{minted}{lean4}
theorem fee_tax_commute (amount fee tax : Money) :
    addFee fee (addFixedTax tax amount)
      = addFixedTax tax (addFee fee amount) := by
  unfold addFee addFixedTax
  exact Nat.add_right_comm amount tax fee
\end{minted}
\caption{\texttt{fee\_tax\_commute}}
\label{lst:ch26-fee-tax-commute}
\end{listing}

この定理は、\texttt{amount}、\texttt{fee}、\texttt{tax} のすべての自然数について成り立ちます。
テストならいくつかの例を試すだけですが、Lean の定理として書くと、モデル化した範囲の全入力について主張していることになります。

図\ref{fig:ch26-commuting-operations}の上段は、加算どうしなので可換です。

下段は \texttt{Nat} の切り捨て減算が途中で \texttt{0} に達するため、二つの経路が一致しない場合を表します。

\FigurePlaceholder{fig-ch26-commuting-operations.png}{0.90}{可換な操作と反例}{fig:ch26-commuting-operations}
\FloatBarrier

\section{条件なしには順序を入れ替えられない処理}

次に、本章の \texttt{Nat} モデルで順序交換が常には成り立たない例を見ます。

たとえば、商品価格が 500 円、割引が 1000 円、手数料が 100 円だとします。
先に割引すると、Lean の自然数モデルでは 0 円になり、その後で手数料を足して 100 円になります。
先に手数料を足すと 600 円になり、そこから 1000 円を引くので 0 円になります。
結果は 100 円と 0 円で一致しません。

\begin{listing}[htbp]
\begin{minted}{lean4}
def discountThenFee (amount discount fee : Money) : Money :=
  addFee fee (applyDiscount discount amount)

def feeThenDiscount (amount discount fee : Money) : Money :=
  applyDiscount discount (addFee fee amount)

theorem discount_fee_order_counterexample :
    discountThenFee 500 1000 100
      ≠ feeThenDiscount 500 1000 100 := by
  decide
\end{minted}
\caption{\texttt{discount\_fee\_order\_counterexample}}
\label{lst:ch26-discount-fee-counterexample}
\end{listing}

ここでは、\texttt{by decide} によって、具体的な自然数計算を Lean に確認させています。
これは「割引と手数料は常に入れ替えられない」という定理ではなく、「入れ替えられるという一般定理は成り立たない」ことを示す反例です。

この不一致の直接の原因は、\texttt{Nat} の減算が \texttt{0} で止まることです。

\texttt{discount <= amount} という条件のもとでは、この固定額割引と固定手数料は順序交換できます。

符号付き整数上の固定額減算でも、同じ反例は生じません。

手数料を割引対象に含めるという業務上の違いを表したい場合は、率や上限を持つ別の割引関数をモデル化する必要があります。

\begin{warningbox}
反例は、仕様レビューで非常に役に立ちます。
「この順序変更は安全ですか」と聞かれたとき、一つでも結果が変わる入力があれば、一般的なリファクタリングとしては安全ではありません。
業務上その順序変更を採用するなら、それはリファクタリングではなく仕様変更として扱うべきです。
\end{warningbox}

\section{等式で証明できる性質}

次に、リファクタリングの安全性を等式で証明します。
標準実装を \texttt{totalV1}、リファクタリング後の実装を \texttt{totalV2} とします。

\begin{listing}[htbp]
\begin{minted}{lean4}
def totalV1 (base discount fee tax : Money) : Money :=
  addFixedTax tax
    (addFee fee (applyDiscount discount base))

def totalV2 (base discount fee tax : Money) : Money :=
  addFee fee
    (addFixedTax tax (applyDiscount discount base))
\end{minted}
\caption{\texttt{totalV1} と \texttt{totalV2}}
\label{lst:ch26-total-v1-v2}
\end{listing}

\texttt{totalV1} は、割引、手数料、固定税の順に処理します。
\texttt{totalV2} は、割引後に固定税を足し、その後に手数料を足します。
つまり、割引の位置は変えていません。
割引後の金額に、固定税と固定手数料を足す順序だけを入れ替えています。

この変更は安全です。
割引後の金額を \texttt{x} と見れば、\texttt{(x + fee) + tax} と \texttt{(x + tax) + fee} が等しいためです。

\begin{listing}[htbp]
\begin{minted}{lean4}
theorem total_refactor_safe
    (base discount fee tax : Money) :
    totalV1 base discount fee tax
      = totalV2 base discount fee tax := by
  unfold totalV1 totalV2 addFixedTax addFee applyDiscount
  exact Nat.add_right_comm (base - discount) fee tax
\end{minted}
\caption{\texttt{total\_refactor\_safe}}
\label{lst:ch26-total-refactor-safe}
\end{listing}

この定理は、\texttt{base}、\texttt{discount}、\texttt{fee}、\texttt{tax} のすべての自然数について成り立ちます。
テストデータを 10 個用意するのではなく、モデル化した全入力に対して、リファクタリング前後の結果が同じであることを確認しています。

同じ証明を、\texttt{calc} を使ってレビューしやすい形で書くこともできます。

\begin{listing}[htbp]
\begin{minted}{lean4}
theorem total_refactor_safe_calc
    (base discount fee tax : Money) :
    totalV1 base discount fee tax
      = totalV2 base discount fee tax := by
  unfold totalV1 totalV2 addFixedTax addFee applyDiscount
  calc
    (base - discount + fee) + tax
        = (base - discount + tax) + fee := by
            exact Nat.add_right_comm (base - discount) fee tax
\end{minted}
\caption{\texttt{total\_refactor\_safe\_calc}}
\label{lst:ch26-total-refactor-calc}
\end{listing}

\texttt{calc} は、式変形のログとして読めます。
料金計算のようにレビューで説明が必要なコードでは、この形の証明が読みやすい場合があります。

\section{不等式で証明すべき性質}

料金計算では、すべてが等式になるわけではありません。
割引は、金額を同じか小さくする操作です。
手数料は、金額を同じか大きくする操作です。
このような性質は、等式ではなく不等式で表します。

\begin{definitionbox}
等式は「結果が同じである」ことを表します。
不等式は「結果がある範囲に収まる」ことを表します。
料金計算では、リファクタリングの安全性は等式で表し、割引後の金額が元の金額を超えないことや、手数料加算後の金額が元の金額を下回らないことは不等式で表します。
\end{definitionbox}

Lean では、次のように書けます。

\begin{listing}[htbp]
\begin{minted}{lean4}
theorem discount_never_increases
    (amount discount : Money) :
    applyDiscount discount amount ≤ amount := by
  unfold applyDiscount
  exact Nat.sub_le amount discount

theorem addFee_never_decreases
    (amount fee : Money) :
    amount ≤ addFee fee amount := by
  unfold addFee
  exact Nat.le_add_right amount fee
\end{minted}
\caption{割引・手数料の上下界}
\label{lst:ch26-inequalities}
\end{listing}

さらに、標準順序の合計が割引直後の金額以上であることも証明できます。

\begin{listing}[htbp]
\begin{minted}{lean4}
theorem total_ge_discounted
    (base discount fee tax : Money) :
    applyDiscount discount base
      ≤ totalV1 base discount fee tax := by
  unfold totalV1 addFixedTax addFee
  exact Nat.le_trans
    (Nat.le_add_right (applyDiscount discount base) fee)
    (Nat.le_add_right
      (applyDiscount discount base + fee) tax)
\end{minted}
\caption{\texttt{total\_ge\_discounted}}
\label{lst:ch26-total-ge-discounted}
\end{listing}

この種の不等式は、料金仕様でよく使います。
たとえば、「割引適用後の金額は元の金額を超えない」「手数料を足した後の金額は減らない」「最低請求額を適用した後の金額は最低額以上である」といった性質です。

\section{丸めは仕様境界として扱う}

最後に、丸めを扱います。
丸めは、料金計算で特に危険な部分です。
丸めは、足し算や掛け算と同じようには振る舞わないためです。

たとえば、10\% の税を考えます。
9 円の明細が二つあるとします。
明細ごとに税を計算して切り捨てるなら、各明細の税は 0 円で、合計税額も 0 円になります。
一方、先に 18 円へ合算してから税を計算すると、税額は 1 円になります。

これは、どちらかが数学的に間違っているという話ではありません。
仕様として「明細ごとに丸める」のか「合計後に丸める」のかを決める必要があるという話です。

\begin{listing}[htbp]
\begin{minted}{lean4}
def taxPerLine2 (a b : Money) : Money :=
  taxOnly 1 10 a + taxOnly 1 10 b

def taxOnSum2 (a b : Money) : Money :=
  taxOnly 1 10 (a + b)

#eval taxPerLine2 9 9  -- => 0
#eval taxOnSum2 9 9  -- => 1

theorem rounding_position_counterexample :
    taxPerLine2 9 9 ≠ taxOnSum2 9 9 := by
  decide
\end{minted}
\caption{\texttt{rounding\_position\_counterexample}}
\label{lst:ch26-rounding-counterexample}
\end{listing}

このコードでは、明細ごとに税を計算する \texttt{taxPerLine2} と、合計してから税を計算する \texttt{taxOnSum2} を分けて定義しています。

\texttt{\#eval} の二行は具体例の結果を観察するためのもので、\texttt{rounding\_position\_counterexample} はその二つが等しくないことを Lean に確認させています。

図\ref{fig:ch26-rounding-boundary}は、各明細で整数除算してから足す経路と、先に合計してから一度だけ整数除算する経路を分けています。

反例は、この二つの経路が一般には可換でないことを示します。

\FigurePlaceholder{fig-ch26-rounding-boundary.png}{0.90}{丸め位置}{fig:ch26-rounding-boundary}

\begin{warningbox}
丸めを含む処理では、「あとでまとめて計算しても同じはず」と考えてはいけません。
丸めの位置は、実装の都合ではなく仕様の一部です。
Lean では反例を小さく示すことで、この順序変更が安全なリファクタリングではないことを確認できます。
\end{warningbox}

\section{圏論的に見る：金額上の処理と可換性}

本章の例には、圏論の大きな理論は必要ありません。
しかし、圏論の言葉で見ると整理しやすくなります。

金額型 \texttt{Money} を対象と見ます。
割引、手数料、税、丸めは、\texttt{Money -> Money} の関数です。
つまり、同じ対象から同じ対象への射です。
このような射は合成できます。

料金計算パイプラインは、\texttt{Money -> Money} の処理の合成です。

順序交換の安全性は、対応する二つの合成が同じ関数になるか、つまり図式が可換になるかで判定できます。

固定手数料と固定税は可換です。
割引と手数料は、一般には可換ではありません。
丸めと合算も、一般には可換ではありません。
これらを区別するだけでも、設計レビューでの会話はかなり明確になります。

\section{仕様レビューでの使い方}

料金計算の Lean モデルを実務で使うときは、次の観点を明示するとよいでしょう。

\begin{table}[htbp]
\centering
\caption{料金計算のレビュー観点}
\label{tab:ch26-review-points}
\begin{tabular}{lll}
\toprule
観点 & Lean で表すもの & レビューで確認すること \\
\midrule
順序 & 関数合成 & 仕様変更かリファクタリングか \\
等価性 & 等式定理 & 全入力で結果が変わらないか \\
範囲保証 & 不等式定理 & 金額が期待範囲に収まるか \\
反例 & \texttt{≠} の定理 & 入れ替え不可であること \\
丸め & 具体例と反例 & 丸め位置が仕様に明記されているか \\
境界 & 注意書き & 実装・外部仕様との差分 \\
\bottomrule
\end{tabular}
\end{table}

Lean で証明したからといって、本番コード全体が正しいとは限りません。
証明したのは、あくまで本章の小さなモデルです。
実務では、次の境界を文書に残す必要があります。

\begin{itemize}
\item 実装の金額型と Lean の \texttt{Money} が対応しているかを確認します。
\item 実装の丸め規則と Lean の整数除算が対応しているかを確認します。
\item 税率、割引上限、無料条件などの外部仕様がどこに定義されているかを確認します。
\item 本章の証明が、仕様変更ではなくリファクタリングの安全性を示しているのかを確認します。
\end{itemize}

\begin{warningbox}
「Lean で料金計算を証明した」と述べるときは、必ず「どのモデルで、どの性質を証明したか」を添えます。
特に、丸め、税法、通貨単位、外部 API、DB に保存された値の整合性は、別途仕様、テスト、監査が必要になります。
\end{warningbox}

\section{本章のまとめ}

料金計算は、税、割引、手数料、丸めの順序によって結果が変わります。

固定税と固定手数料のように、足し算だけで表せる処理は、順序交換を等式で証明できます。

割引と手数料、丸めと合算は、一般には順序を入れ替えられません。
反例を Lean で示せます。

リファクタリングの安全性は、変更前後の関数が同じ結果を返すという等式として書けます。

割引が金額を増やさないことや、手数料が金額を減らさないことは、不等式で書きます。

Lean の証明は、仕様の核を検査する道具であり、税法や外部システムを含む本番全体の保証ではありません。

\section{次章への接続}

次章からは、仕様証明の考え方をマイグレーションへ広げます。
料金計算では、処理順序と意味保存を確認しました。
マイグレーションでは、データ形式を変換したときに情報が保存されるか、元へ戻せるか、どの条件なら安全かを扱います。
ここで使った「証明できる性質を小さく切り出す」という考え方は、そのまま次の Part でも使います。
